\chapter{Standaardfuncties}\label{ch:b1:functions}

De analyse is niet nuttiger dan de voorraad functies die je beheerst. Aan
de \omterm{def:b1:logic:sets}{verzameling} die uit het bovenbouwvolume is overgeërfd --- machten,
exponentiële functie, logaritme, goniometrische functies --- voegt dit
hoofdstuk hun inverse functies ($\arcsin$, $\arccos$, $\arctan$) en de
hyperbolische familie toe. Afgeleiden worden vrijelijk gebruikt zoals ze
in het bovenbouwvolume voorkwamen; de theorie achter de inverse functies
wordt vervolledigd in \cref{ch:b1:continuity,ch:b1:derivative}.

\section{Exponentiële functie, logaritme, machten}

\begin{proposition}[Herhaling en karakterisering]\label{prop:b1:functions:expln}
$\exp \colon \R \to \intoo{0}{+\infty}$ en $\ln \colon
\intoo{0}{+\infty} \to \R$ zijn elkaars inverse bijecties, strikt
stijgend, met
\[
\exp(x + y) = \exp x \exp y, \qquad \ln(xy) = \ln x + \ln y,
\qquad \exp' = \exp, \qquad \ln'(x) = \frac 1x .
\]
Bovendien is $\exp$ de \emph{enige} afleidbare functie $f \colon \R \to
\R$ met $f' = f$ en $f(0) = 1$.
\end{proposition}

\begin{proof}
De herhaling is stof uit het bovenbouwvolume. Voor de uniciteit: zij $f'
= f$ en $f(0) = 1$, en zet $g(x) = f(x)\,\eu^{-x}$. Dan is $g' =
f'\eu^{-x} - f\eu^{-x} = 0$, zodat $g$ constant gelijk aan $g(0) = 1$
is: $f = \exp$.
\end{proof}

\begin{definition}[Algemene machten]\label{def:b1:functions:powers}
Voor $x > 0$ en $\alpha \in \R$ zet men $\;x^\alpha = \eu^{\alpha \ln
x}$\index{machtsfunctie}. Voor $a > 0$ met $a \neq 1$ is de
\emph{logaritme met grondtal $a$} de functie $\log_a x = \frac{\ln
x}{\ln a}$, de inverse van $x \mapsto a^x$.
\end{definition}

\begin{example}[Exponentiële vergelijkingen oplossen]\label{ex:b1:functions:expequations}
Los $2^x = 5^{\,x-1}$ op in $\R$. Beide leden zijn positief, dus neem
logaritmen --- een omkeerbare stap:
\[
x\ln2 = (x - 1)\ln5
\iff x(\ln2 - \ln5) = -\ln5
\iff x = \frac{\ln5}{\ln5 - \ln2} = \frac{\ln 5}{\ln\frac52}
\approx 1.756 .
\]
Los vervolgens $x^{\sqrt2} = 3$ op voor $x > 0$: verhef tot de macht
$\frac1{\sqrt2}$ (dat wil zeggen, pas de inverse bijectie toe): $x =
3^{1/\sqrt2} = \eu^{(\ln 3)/\sqrt2} \approx 2.175$. Het inzicht: elke
vergelijking waarin machten door elkaar lopen ontrolt zich via $\ln$ en
$\exp$, omdat de definitie $x^\alpha = \eu^{\alpha\ln x}$ elk
machtsgereken tot rekenen met exponenten herleidt --- maar alleen op
het domein $x > 0$ waar die definitie leeft.
\end{example}

\begin{example}[Verdubbelingstijden]\label{ex:b1:functions:doubling}
Een grootheid groeit met $3\%$ per stap: na $n$ stappen is ze
vermenigvuldigd met $(1.03)^n$. Wanneer verdubbelt ze? Los $(1.03)^n
\geq 2$ op:
\[
n\ln(1.03) \geq \ln2
\iff n \geq \frac{\ln 2}{\ln 1.03}
= \frac{0.6931}{0.02956} \approx 23.4 ,
\]
zodat de eerste verdubbeling bij stap $24$ optreedt. (De ``regel van
$72$'' van de financiers, die de verdubbelingstijd schat als $72$
gedeeld door het percentage, is deze berekening met de benadering
$\ln(1 + x) \approx x$, gekwantificeerd in \cref{ch:b1:taylor}.) Over
exponentiële processen redeneer je het best via hun logaritmen: op die
schaal is de groei lineair en worden vragen delingen.
\end{example}

\begin{example}[Hoe lang is $2^{2026}$?]\label{ex:b1:functions:digits}
Het aantal decimale cijfers van een geheel getal $N \geq 1$ is
$\floor{\log_{10} N} + 1$ (immers, $N$ heeft $d$ cijfers precies
wanneer $10^{d-1} \leq N < 10^d$, oftewel $d - 1 \leq \log_{10}N < d$).
Voor $N = 2^{2026}$:
\[
\log_{10} 2^{2026} = 2026\,\log_{10}2
= 2026 \times 0.301030 = 609.887 ,
\]
dus $2^{2026}$ heeft $610$ cijfers. Het decimale deel levert een bonus:
$10^{0.887} \approx 7.7$, zodat het getal met een $7$ \emph{begint}.
Eén vermenigvuldiging beantwoordde een vraag over een getal dat niemand
ooit zal uitschrijven --- logaritmen persen multiplicatieve grootte
samen tot additieve, en dat is hun hele historische bestaansreden
(\cref{ex:b1:structures:expmorphism} maakt die zin precies).
\end{example}

\begin{proposition}[Machtsregels en groeivergelijking]\label{prop:b1:functions:powerrules}
Voor $x, y > 0$ en $\alpha, \beta \in \R$:
\[
x^{\alpha+\beta} = x^\alpha x^\beta, \quad
(x^\alpha)^\beta = x^{\alpha\beta}, \quad
(xy)^\alpha = x^\alpha y^\alpha, \quad
(x^\alpha)' = \alpha\, x^{\alpha - 1}.
\]
Groeischaal\index{groeivergelijking} als $x \to +\infty$, voor elke
$\alpha > 0$ en $\beta > 0$:
\[
\frac{(\ln x)^\beta}{x^\alpha} \longrightarrow 0,
\qquad
\frac{x^\alpha}{\eu^{\beta x}} \longrightarrow 0 .
\]
\end{proposition}

\begin{proof}
De identiteiten schrijven die van $\exp$ en $\ln$ over via de definitie.
In detail voor de tweede (de minst voor de hand liggende): $x^\alpha > 0$
en $\ln(x^\alpha) = \alpha\ln x$ (pas $\ln$ toe op de definitie), zodat
\[
(x^\alpha)^\beta = \eu^{\beta\ln(x^\alpha)}
= \eu^{\beta\alpha\ln x} = x^{\alpha\beta} ;
\]
de eerste en de derde zijn dezelfde ontrollingen van twee regels via
$\exp(u + v) = \exp u\exp v$ en $\ln(xy) = \ln x + \ln y$. De afgeleide
is de kettingregel: $(\eu^{\alpha\ln x})' = \frac{\alpha}{x}
\eu^{\alpha\ln x}$.

Vergelijkingen: uit $\frac{\ln t}{t} \to 0$ (bewezen in het
bovenbouwvolume) substitueer je $t = x^{\alpha/\beta}$, wat $\frac{\ln
x}{x^{\alpha/\beta}} \to 0$ geeft, en verhef je tot de macht $\beta$.
Voor de tweede substitueer je $x = \ln u$ in de eerste.
\end{proof}

\begin{example}[Twee limieten die elke lezer moet bezitten]\label{ex:b1:functions:xx}
Wat zijn $\lim_{x \to 0^+} x^x$ en $\lim_{x \to +\infty} x^{1/x}$?
Beide machten zijn via de exponentiële functie gedefinieerd, dus
beslecht je eerst de exponent:
\[
x^x = \eu^{x\ln x}, \qquad
x\ln x = -\frac{\ln(1/x)}{1/x} \longrightarrow 0
\quad (x \to 0^+),
\qquad\text{dus } x^x \longrightarrow \eu^0 = 1 ;
\]
en $x^{1/x} = \eu^{(\ln x)/x} \to \eu^0 = 1$ als $x \to +\infty$,
rechtstreeks door de \omterm{prop:b1:functions:powerrules}{groeivergelijking}. Het inzicht: een onbepaalde
macht (van de vorm $0^0$ of $\infty^0$) pak je \emph{altijd} aan door
$u^v = \eu^{v\ln u}$ te schrijven en het product $v \ln u$ te
analyseren --- nooit door base en exponent afzonderlijk te raden. De
functie $x \mapsto \frac{\ln x}{x}$, die beide limieten besliste, wordt
uitputtend bestudeerd in de weekendopgave van dit hoofdstuk.
\end{example}

\begin{omfigure}
\begin{tikzpicture}
\begin{axis}[omaxis, width=9.5cm, height=6cm,
    xmin=0, xmax=5.2, ymin=-1.8, ymax=8.5,
    xtick={1,2,3,4,5}, ytick={2,4,6,8}]
  \addplot[omMeth, thick, domain=0.02:5, samples=200] {ln(x)}
    node[pos=0.97, below, font=\small] {$\ln x$};
  \addplot[omDef, thick, domain=0:5, samples=100] {sqrt(x)}
    node[pos=0.9, above, font=\small] {$\sqrt x$};
  \addplot[omProp, thick, domain=0:5] {x}
    node[pos=0.72, left, font=\small] {$x$};
  \addplot[omThm, thick, domain=0:2.14, samples=120] {exp(x)}
    node[pos=0.85, left, font=\small] {$\eu^x$};
\end{axis}
\end{tikzpicture}

{\small De groeischaal van \cref{prop:b1:functions:powerrules}: bij de
rechterrand is $\ln x$ amper voorbij $1.6$, terwijl $\eu^x$ het kader
al verlaten heeft. Elke verhouding (logaritme boven macht, macht boven
exponentiële) gaat naar $0$ --- de tekening suggereert alleen wat de
substituties in het bewijs exact maken.}
\end{omfigure}

\begin{remark}[Veelgemaakte fouten met machten en logaritmen]\label{rem:b1:functions:pitfalls}
\begin{enumerate}
  \item \emph{Domeinen.} $x^\alpha$ voor irrationale $\alpha$ vergt $x
        > 0$; $(-8)^{1/3}$ vermijd je beter ten gunste van ``de reële
        derdemachtswortel van $-8$'', want de regel $(x^\alpha)^\beta =
        x^{\alpha\beta}$ faalt stilzwijgend op negatieve getallen:
        $\bigl((-8)^2\bigr)^{1/6} = 2 \neq -2$.
  \item \emph{$\sqrt{x^2} = \abs x$}, niet $x$: de absolute waarde
        vergeten is de klassieke bron van zoekgeraakte negatieve
        oplossingen.
  \item \emph{$\ln(xy) = \ln x + \ln y$ vergt $x, y > 0$.} Voor
        negatieve argumenten kan $\ln(xy)$ gedefinieerd zijn terwijl
        het rechterlid dat niet is.
  \item \emph{Welke functie groeit?} In $x^\alpha$ varieert het
        \emph{grondtal}; in $a^x$ de \emph{exponent}. Hun groei
        verschilt hemelsbreed (\cref{prop:b1:functions:powerrules}), en
        hybride uitdrukkingen als $x^{\ln x}$ of $x^{1/x}$ moeten als
        $\eu^{(\cdot)}$ herschreven worden vóór elke redenering ---
        zoals in \cref{ex:b1:functions:xx}.
\end{enumerate}
\end{remark}

\section{Inverse goniometrische functies}

\begin{definition}[$\arcsin$, $\arccos$, $\arctan$]\label{def:b1:functions:arc}
De beperkingen
\[
\sin \colon \intcc{-\tfrac\pi2}{\tfrac\pi2} \to \intcc{-1}{1},
\qquad
\cos \colon \intcc{0}{\pi} \to \intcc{-1}{1},
\qquad
\tan \colon \intoo{-\tfrac\pi2}{\tfrac\pi2} \to \R
\]
zijn strikt monotone bijecties. Hun inverse \omterm{def:b1:logic:map}{afbeeldingen} worden
$\arcsin$\index{boogsinus}, $\arccos$\index{boogcosinus} en
$\arctan$\index{boogtangens} genoteerd. Zo is bijvoorbeeld $y = \arcsin
x$ (met $x \in \intcc{-1}{1}$) dé hoek in
$\intcc{-\frac\pi2}{\frac\pi2}$ waarvan de sinus $x$ is.
\end{definition}

\begin{omfigure}
\begin{tikzpicture}
\begin{axis}[omaxis, width=5.6cm, height=5.6cm,
    xmin=-1.25, xmax=1.25, ymin=-1.8, ymax=1.8,
    xtick={-1,1}, ytick={-1.5708,1.5708},
    yticklabels={$-\frac{\pi}{2}$,$\frac{\pi}{2}$}]
  \addplot[omDef, very thick, domain=-1:1, samples=200]
    {rad(asin(x))};
\end{axis}
\begin{scope}[xshift=5.3cm]
\begin{axis}[omaxis, width=5.6cm, height=5.6cm,
    xmin=-1.25, xmax=1.25, ymin=-0.4, ymax=3.4,
    xtick={-1,1}, ytick={1.5708,3.1416},
    yticklabels={$\frac{\pi}{2}$,$\pi$}]
  \addplot[omThm, very thick, domain=-1:1, samples=200]
    {rad(acos(x))};
\end{axis}
\end{scope}
\begin{scope}[xshift=10.6cm]
\begin{axis}[omaxis, width=5.6cm, height=5.6cm,
    xmin=-4, xmax=4, ymin=-1.9, ymax=1.9,
    xtick={-2,2}, ytick={-1.5708,1.5708},
    yticklabels={$-\frac{\pi}{2}$,$\frac{\pi}{2}$}]
  \addplot[omMeth, very thick, domain=-4:4, samples=200]
    {rad(atan(x))};
  \addplot[dashed, gray, domain=-4:4] {1.5708};
  \addplot[dashed, gray, domain=-4:4] {-1.5708};
\end{axis}
\end{scope}
\end{tikzpicture}

{\small Van links naar rechts: $\arcsin$, $\arccos$, $\arctan$. Elk
erft zijn grafiek van de beperkte directe functie door spiegeling in de
rechte $y = x$.}
\end{omfigure}

\begin{proposition}[Afgeleiden]\label{prop:b1:functions:arcderiv}
Op het inwendige van hun domeinen geldt
\[
\arcsin' x = \frac{1}{\sqrt{1 - x^2}},
\qquad
\arccos' x = \frac{-1}{\sqrt{1 - x^2}},
\qquad
\arctan' x = \frac{1}{1 + x^2}.
\]
\end{proposition}

\begin{proof}
We lopen op de regel voor de afgeleide van een inverse functie vooruit,
die in \cref{ch:b1:derivative} bewezen wordt: is $f$ een afleidbare
bijectie met $f' \neq 0$, dan is $(f^{-1})'(x) = \frac{1}{f'(f^{-1}(x))}$.
Voor $\arcsin$ geeft dat, met $y = \arcsin x$,
\[
\arcsin' x = \frac{1}{\cos y}
= \frac{1}{\sqrt{1 - \sin^2 y}} = \frac{1}{\sqrt{1 - x^2}},
\]
waarbij $\cos y = +\sqrt{1 - \sin^2 y}$ omdat $y \in
\intoo{-\frac\pi2}{\frac\pi2}$ afdwingt dat $\cos y > 0$. Analoog is
$\arccos' x = \frac{1}{-\sin y}$ met $\sin y > 0$ op $\intoo{0}{\pi}$,
en $\arctan' x = \frac{1}{1 + \tan^2 y} = \frac{1}{1 + x^2}$ met $\tan'
= 1 + \tan^2$.
\end{proof}

\begin{example}[Een verborgen constante herkennen]\label{ex:b1:functions:hiddenconst}
Bestudeer $g(x) = \arctan\dfrac{1 - x}{1 + x}$ op $\intoo{-1}
{+\infty}$. De kettingregel en een korte berekening geven
\[
g'(x) = \frac{1}{1 + \bigl(\frac{1-x}{1+x}\bigr)^2}\cdot
\frac{-(1+x) - (1-x)}{(1+x)^2}
= \frac{-2}{(1+x)^2 + (1-x)^2}
= \frac{-1}{1 + x^2} ,
\]
want $(1+x)^2 + (1-x)^2 = 2 + 2x^2$. Dus $g' = -\arctan'$: de functie
$g + \arctan$ heeft afgeleide nul op het \emph{interval}
$\intoo{-1}{+\infty}$ en is daar dus constant; haar waarde in $x = 0$ is
$\arctan 1 + \arctan 0 = \frac\pi4$. Besluit:
\[
\arctan\frac{1 - x}{1 + x} = \frac\pi4 - \arctan x
\qquad (x > -1) .
\]
Op $\intoo{-\infty}{-1}$ geldt dezelfde berekening van de afgeleide,
maar is de constante een andere ($-\frac{3\pi}4$: bereken de limiet
voor $x \to -\infty$). Het inzicht: ``afgeleide nul, dus constant'' is
een \omterm{def:b1:logic:statement}{uitspraak} per interval --- precies de subtiliteit die in
\cref{exo:b1:functions:6} en \cref{exo:b1:functions:10} uitgebuit
wordt.
\end{example}

\begin{proposition}[Standaardidentiteiten]\label{prop:b1:functions:arcidentities}
\begin{enumerate}
  \item Voor $x \in \intcc{-1}{1}$: $\arcsin x + \arccos x =
        \dfrac{\pi}{2}$.
  \item Voor $x > 0$: $\arctan x + \arctan\dfrac 1x = \dfrac{\pi}{2}$
        (en $-\frac\pi2$ voor $x < 0$).
  \item $\sin(\arccos x) = \cos(\arcsin x) = \sqrt{1 - x^2}$;
        $\;\tan(\arcsin x) = \dfrac{x}{\sqrt{1 - x^2}}$ voor
        $\abs x < 1$.
\end{enumerate}
\end{proposition}

\begin{proof}
(1) De afgeleide van $x \mapsto \arcsin x + \arccos x$ is nul op
$\intoo{-1}{1}$ (\cref{prop:b1:functions:arcderiv}), zodat de functie
daar constant is en gelijk aan haar waarde $\frac\pi2$ in $0$; de
randwaarden $\pm 1$ ga je rechtstreeks na ($\frac\pi2 + 0$ en
$-\frac\pi2 + \pi$).

(2) Dezelfde methode op $\intoo{0}{+\infty}$: de afgeleide is
$\frac{1}{1+x^2} + \frac{-1/x^2}{1 + 1/x^2} = \frac{1}{1+x^2} -
\frac{1}{x^2 + 1} = 0$, en in $x = 1$ is de som $2\arctan 1 =
\frac\pi2$. Voor $x < 0$ gebruik je dat $\arctan$ oneven is.

(3) Met $y = \arccos x \in \intcc{0}{\pi}$ is $\sin y \geq 0$, dus $\sin
y = \sqrt{1 - \cos^2 y} = \sqrt{1 - x^2}$; analoog voor $\cos(\arcsin
x)$, en de formule voor de tangens is het quotiënt.
\end{proof}

\begin{remark}
$\arcsin(\sin\theta) = \theta$ geldt \emph{alleen} voor $\theta \in
\intcc{-\frac\pi2}{\frac\pi2}$: zo is $\arcsin(\sin\pi) = 0 \neq \pi$.
De samenstelling in de andere richting, $\sin(\arcsin x) = x$, is geldig
op heel $\intcc{-1}{1}$.
\end{remark}

\begin{example}[Terugvouwen naar het hoofdinterval]
\label{ex:b1:functions:wrap}
Bereken
\[
\arctan\Bigl(\tan\frac{3\pi}4\Bigr)
\qquad\text{en}\qquad
\arctan\Bigl(\tan\frac{17\pi}5\Bigr).
\]
Het recept: vervang de hoek door de unieke hoek uit
$\intoo{-\frac\pi2}{\frac\pi2}$ met dezelfde tangens, dat wil zeggen
trek het juiste veelvoud van $\pi$ af (de periode van $\tan$). Eerst:
$\frac{3\pi}4 - \pi = -\frac\pi4$, dus het antwoord is $-\frac\pi4$.
Dan: $\frac{17\pi}5 - 3\pi = \frac{2\pi}5 \in
\intoo{-\frac\pi2}{\frac\pi2}$, dus het antwoord is $\frac{2\pi}5$. De
berekening is een vermomde euclidische deling van de hoek door $\pi$
--- en de analoge recepten voor $\arcsin$ (vouw terug naar
$\intcc{-\frac\pi2}{\frac\pi2}$, periode $2\pi$) en $\arccos$ (vouw
terug naar $\intcc0\pi$) sturen \cref{exo:b1:functions:1} en het
stuksgewijze antwoord van \cref{exo:b1:functions:10}.
\end{example}

\begin{example}[Boogtangenten veilig optellen]\label{ex:b1:functions:arctansum}
We bewijzen
\[
\arctan\frac12 + \arctan\frac15 + \arctan\frac18 = \frac\pi4 .
\]
Twee ingrediënten: de somformule voor de tangens en --- de stap die
beginners vergeten --- een \emph{lokalisering} van de som. Ten eerste
geeft $\tan(u + v) = \frac{\tan u + \tan v}{1 - \tan u\tan v}$ met $u =
\arctan\frac12$ en $v = \arctan\frac15$
\[
\tan(u + v) = \frac{\frac12 + \frac15}{1 - \frac1{10}}
= \frac{7/10}{9/10} = \frac79,
\]
en vervolgens
\[
\tan\Bigl(u + v + \arctan\frac18\Bigr)
= \frac{\frac79 + \frac18}{1 - \frac7{72}}
= \frac{65/72}{65/72} = 1 .
\]
Ten tweede ligt elk van de drie hoeken in $\intoo0{\frac\pi4}$ (hun
argumenten zijn kleiner dan $1$), zodat de som in
$\intoo0{\frac{3\pi}4}$ ligt; de enige hoek daar met tangens $1$ is
$\frac\pi4$. Zonder de lokalisering zou de conclusie ``$\frac\pi4$ op
een veelvoud van $\pi$ na'' luiden --- een half bewijs. Dezelfde
tweestapsdiscipline stuurt \cref{exo:b1:functions:8} en
\cref{exo:b1:functions:12}.
\end{example}

\begin{example}[Een vergelijking met boogtangenten]\label{ex:b1:functions:arctaneq}
Los $\arctan x + \arctan 2x = \dfrac\pi4$ op. Lokaliseer eerst: het
linkerlid heeft het teken van $x$ (beide termen hebben dat), dus heeft
elke oplossing $x > 0$, en dan ligt de som in $\intoo0\pi$. Neem
tangenten (\omterm{def:b1:logic:inj}{injectief} op geen enkel interval van lengte $\pi$, maar in
combinatie met de lokalisering volstaat dit): de somformule geeft
\[
\tan\bigl(\arctan x + \arctan 2x\bigr)
= \frac{3x}{1 - 2x^2} = 1
\iff 2x^2 + 3x - 1 = 0
\iff x = \frac{-3 \pm \sqrt{17}}4 .
\]
De negatieve wortel valt af door de lokalisering; voor de positieve
kandidaat $x_0 = \frac{\sqrt{17} - 3}4 \approx 0.28$ ligt de som in
$\intoo0\pi$ met tangens $1$, en de enige zulke hoek is $\frac\pi4$ (op
$\intoo{\frac\pi2}\pi$ is de tangens negatief): $x_0$ is de unieke
oplossing. Let op de vorm van het argument: tangenten nemen kan
oplossingen creëren, nooit doen verdwijnen, dus los je eerst de
veeltermvergelijking op \emph{en} filter je daarna met de lokalisering
--- dezelfde controlediscipline achteraf als bij het kwadrateren van
een vergelijking.
\end{example}

\section{Hyperbolische functies}

\begin{definition}[$\cosh$, $\sinh$, $\tanh$]\label{def:b1:functions:hyperbolic}
Voor $x \in \R$:
\[
\cosh x = \frac{\eu^x + \eu^{-x}}{2},
\qquad
\sinh x = \frac{\eu^x - \eu^{-x}}{2},
\qquad
\tanh x = \frac{\sinh x}{\cosh x}
\]
(de \emph{hyperbolische functies}\index{hyperbolische functies}: cosinus,
sinus en tangens hyperbolicus). $\cosh$ is even, $\sinh$ en $\tanh$ zijn
oneven.
\end{definition}

\begin{proposition}[Basiseigenschappen]\label{prop:b1:functions:hyprules}
Voor alle $x, y \in \R$:
\begin{enumerate}
  \item $\cosh^2 x - \sinh^2 x = 1$;
  \item $\cosh' = \sinh$, $\;\sinh' = \cosh$, $\;\tanh' = 1 - \tanh^2
        = \dfrac{1}{\cosh^2}$;
  \item $\cosh(x+y) = \cosh x\cosh y + \sinh x \sinh y$ en
        $\sinh(x+y) = \sinh x \cosh y + \cosh x \sinh y$;
  \item $\sinh$ is een strikt stijgende bijectie van $\R$ op $\R$;
        $\cosh$ beperkt tot $\R_+$ is een strikt stijgende bijectie op
        $\intco{1}{+\infty}$; $\tanh$ is een strikt stijgende bijectie
        van $\R$ op $\intoo{-1}{1}$.
\end{enumerate}
\end{proposition}

\begin{proof}
(1) $\bigl(\frac{\eu^x + \eu^{-x}}{2}\bigr)^2 - \bigl(\frac{\eu^x -
\eu^{-x}}{2}\bigr)^2 = \frac{(\eu^{2x} + 2 + \eu^{-2x}) - (\eu^{2x} - 2 +
\eu^{-2x})}{4} = 1$.

(2) Differentieer de definiërende formules; voor $\tanh$ geeft de
quotiëntregel $\frac{\cosh^2 - \sinh^2}{\cosh^2}$, wat volgens (1) zowel
$1 - \tanh^2$ als $\frac{1}{\cosh^2}$ is.

(3) Werk de rechterleden uit met de definities. In detail voor de eerste
formule:
\[
\cosh x\cosh y + \sinh x\sinh y
= \frac{(\eu^x + \eu^{-x})(\eu^y + \eu^{-y})
+ (\eu^x - \eu^{-x})(\eu^y - \eu^{-y})}{4} ;
\]
van de acht producten vallen de vier ``gemengde'' ($\eu^{x-y}$,
$\eu^{y-x}$) twee aan twee weg, terwijl $\eu^{x+y}$ en $\eu^{-(x+y)}$
elk tweemaal voorkomen: het totaal is $\frac{2\eu^{x+y} +
2\eu^{-(x+y)}}{4} = \cosh(x+y)$. De formule voor de sinus verloopt
identiek, met de gemengde termen die nu juist overleven.

(4) $\sinh' = \cosh \geq 1 > 0$, dus $\sinh$ is strikt stijgend, met
limieten $\pm\infty$ (dominante term $\pm\frac{\eu^{\abs x}}{2}$); de
\omterm{def:b1:logic:statement}{uitspraak} over de bijectie volgt dan uit de tussenwaardestelling (hier
gebruikt zoals bekend uit het bovenbouwvolume; systematische behandeling
in \cref{ch:b1:continuity}). Op $\R_+$ is $\cosh' = \sinh \geq 0$, alleen
nul in $0$: strikt stijgend van $\cosh 0 = 1$ naar $+\infty$. En $\tanh'
> 0$ met limieten $\pm 1$ in $\pm\infty$; in detail:
\[
\tanh x = \frac{\eu^x - \eu^{-x}}{\eu^x + \eu^{-x}}
= \frac{1 - \eu^{-2x}}{1 + \eu^{-2x}}
\xrightarrow[x \to +\infty]{} 1 ,
\]
na deling van teller en noemer door $\eu^x$, en wegens de oneven
pariteit is de limiet $-1$ in $-\infty$; de strikt stijgende functie
$\tanh$ beeldt $\R$ dus af op $\intoo{-1}1$.
\end{proof}

\begin{omfigure}
\begin{tikzpicture}
\begin{axis}[omaxis, width=6.4cm, height=6cm,
    xmin=-2.6, xmax=2.6, ymin=-3.6, ymax=3.9,
    xtick={-2,-1,1,2}, ytick={-3,-2,-1,1,2,3}]
  \addplot[omDef, very thick, domain=-2.05:2.05, samples=120]
    {cosh(x)} node[pos=0.1, right, font=\small] {$\cosh$};
  \addplot[omThm, very thick, domain=-2.05:2.05, samples=120]
    {sinh(x)} node[pos=0.93, above left, font=\small] {$\sinh$};
  \addplot[dashed, gray, domain=0:2.5] {exp(x)/2};
\end{axis}
\begin{scope}[xshift=7.4cm]
\begin{axis}[omaxis, width=6.4cm, height=6cm,
    xmin=-3.2, xmax=3.2, ymin=-1.5, ymax=1.5,
    xtick={-2,-1,1,2}, ytick={-1,1}]
  \addplot[omMeth, very thick, domain=-3:3, samples=150]
    {tanh(x)} node[pos=0.9, above left, font=\small] {$\tanh$};
  \addplot[dashed, gray, domain=-3:3] {1};
  \addplot[dashed, gray, domain=-3:3] {-1};
\end{axis}
\end{scope}
\end{tikzpicture}

{\small Links: $\cosh$ en $\sinh$, asymptotisch aan $\frac{\eu^x}{2}$
geplakt (streepjeslijn). Rechts: $\tanh$, stijgend van $-1$ naar $1$.}
\end{omfigure}

\begin{remark}[Waarom ``hyperbolisch''?]
Het punt $(\cosh t, \sinh t)$ doorloopt de tak $x > 0$ van de hyperbool
$x^2 - y^2 = 1$ (wegens
\cref{prop:b1:functions:hyprules}~(1)), precies zoals $(\cos t, \sin
t)$ de cirkel $x^2 + y^2 = 1$ doorloopt. Elke goniometrische identiteit
heeft een hyperbolische tegenhanger, met tekenwisselingen die door
$\sinh^2 \leftrightarrow -\sin^2$ geregeld worden.
\end{remark}

\begin{example}[De somformule voor $\tanh$]\label{ex:b1:functions:tanhadd}
Delen we de twee somformules van
\cref{prop:b1:functions:hyprules}~(3) door $\cosh x\cosh y$, dan volgt
\[
\tanh(x + y)
= \frac{\sinh x\cosh y + \cosh x\sinh y}
{\cosh x\cosh y + \sinh x\sinh y}
= \frac{\tanh x + \tanh y}{1 + \tanh x\,\tanh y} ,
\]
de hyperbolische tegenhanger van de somformule voor de tangens --- met
een $+$ waar de goniometrie een $-$ heeft. Een dividend: omdat
$\abs{\tanh} < 1$, is het rechterlid een ``optelregel voor snelheden''
die $\intoo{-1}1$ nooit verlaat: liggen $u, v$ in $\intoo{-1}1$, dan
ligt ook $\frac{u + v}{1 + uv}$ in $\intoo{-1}1$ (schrijf $u = \tanh a$
en $v = \tanh b$, wat kan wegens de bijectiviteit, en lees de formule
achterstevoren). Dat algebraïsch nagaan, zonder \omterm{def:b1:functions:hyperbolic}{hyperbolische functies},
is een licht pijnlijke oefening; parametriseren met $\tanh$ maakt er één
regel van --- dezelfde strategie die de goniometrische functies voor de
eenheidscirkel bieden.
\end{example}

\begin{proposition}[Inverse hyperbolische functies]\label{prop:b1:functions:invhyp}
De inversen die \cref{prop:b1:functions:hyprules}~(4) oplevert hebben
gesloten uitdrukkingen:
\[
\operatorname{arsinh} x = \ln\bigl(x + \sqrt{x^2 + 1}\bigr)
\ (x \in \R),
\qquad
\operatorname{arcosh} x = \ln\bigl(x + \sqrt{x^2 - 1}\bigr)
\ (x \geq 1),
\]
\[
\operatorname{artanh} x = \frac 12 \ln\frac{1 + x}{1 - x}
\ (\abs x < 1),
\]
met respectievelijk de afgeleiden $\frac{1}{\sqrt{x^2+1}}$,
$\frac{1}{\sqrt{x^2-1}}$ ($x > 1$) en $\frac{1}{1 - x^2}$.
\end{proposition}

\begin{proof}
Voor $\operatorname{arsinh}$: los $x = \sinh y = \frac{\eu^y -
\eu^{-y}}{2}$ op. Met $u = \eu^y > 0$ wordt dat $u^2 - 2xu - 1 = 0$, dus
$u = x + \sqrt{x^2 + 1}$ (de wortel $x - \sqrt{x^2+1}$ is negatief) en $y
= \ln(x + \sqrt{x^2+1})$. De twee andere zijn identieke berekeningen
($u^2 - 2xu + 1 = 0$ voor $\operatorname{arcosh}$, met behoud van de
wortel $\geq 1$; een herschikking van twee regels voor
$\operatorname{artanh}$). Afgeleiden: differentieer de logaritmische
uitdrukkingen, bijvoorbeeld
\[
\bigl(\ln(x + \sqrt{x^2+1})\bigr)'
= \frac{1 + \frac{x}{\sqrt{x^2+1}}}{x + \sqrt{x^2+1}}
= \frac{1}{\sqrt{x^2 + 1}} . \qedhere
\]
\end{proof}

\begin{example}[Gesloten vormen aan het werk]\label{ex:b1:functions:arcoshtwo}
De oplossing van $\cosh t = 2$ met $t \geq 0$ is volgens de gesloten
vorm $t = \operatorname{arcosh} 2 = \ln(2 + \sqrt3) \approx 1.317$; de
andere oplossing is $-t$, wegens de even pariteit --- en inderdaad is
$\ln(2 - \sqrt3) = \ln\frac1{2 + \sqrt3} = -\ln(2 + \sqrt3)$: de twee
wortels $u = \eu^t$ van de vierkantsvergelijking $u^2 - 4u + 1 = 0$ zijn
elkaars omgekeerde, zoals hun product $1$ (Vieta) eist. Dit kleine
rekenwerk toont het algemene patroon: hyperbolische vergelijkingen gaan
over in vierkantsvergelijkingen in $\eu^t$, en de symmetrie $t \mapsto
-t$ verschijnt als de symmetrie $u \mapsto 1/u$ van die vergelijking ---
het onthouden waard bij het oplossen van \cref{exo:b1:functions:7}.
\end{example}

\begin{method}[De juiste primitieve vorm kiezen]\label{met:b1:functions:primitive}
De drie patronen van afgeleiden die je voor het integreren
(\cref{ch:b1:integration}) uit het hoofd kent:
\[
\int \frac{\dd x}{1 + x^2} = \arctan x + C,
\qquad
\int \frac{\dd x}{\sqrt{1 - x^2}} = \arcsin x + C,
\qquad
\int \frac{\dd x}{\sqrt{x^2 + 1}} = \operatorname{arsinh} x + C,
\]
en voor een algemene tweedegraadsuitdrukking herleid je tot deze door
het kwadraat af te splitsen en te herschalen.
\end{method}

\begin{remark}[Waar dit hoofdstuk gebruikt wordt]\label{rem:b1:functions:whereused}
Dit hoofdstuk levert het werkvocabulaire van alle analyse die volgt. De
groeischaal van \cref{prop:b1:functions:powerrules} beslist
convergentievragen doorheen \cref{ch:b1:seq,ch:b1:series}; de formules
voor de afgeleiden uit \cref{prop:b1:functions:arcderiv} en
\cref{prop:b1:functions:invhyp} zijn de primitieven die in
\cref{ch:b1:integration} het vaakst nodig zijn, via
\cref{met:b1:functions:primitive}; \omterm{def:b1:functions:hyperbolic}{hyperbolische functies}
parametriseren de oplossingen van de vergelijking $y'' = y$ in
\cref{ch:b1:diffeq} precies zoals de goniometrische functies die van
$y'' = -y$ parametriseren. De karakterisering van $\exp$ door $f' = f$
en $f(0) = 1$ (\cref{prop:b1:functions:expln}) is de eendimensionale
kiem van de theorie van de lineaire differentiaalvergelijkingen, en de
kettinglijn $y = \cosh x$ keert terug bij de vlakke krommen van
\cref{ch:b1:curves}.
\end{remark}

\begin{remark}[Tussenspel: het programma van de inverse functies]\label{rem:b1:functions:interlude}
Dit hoofdstuk voerde viermaal hetzelfde programma uit: beperk een
functie tot ze een strikt monotone bijectie is, geef de inverse een
naam, en vervoer elke formule daardoorheen. Wat bij elke stap
\emph{gebruikt} werd --- dat een continue strikt monotone functie op een
interval een bijectie op een interval is, en dat haar inverse continu is
en vervolgens afleidbaar buiten de kritieke punten --- werd op krediet
van het bovenbouwvolume geleend. Die schuld wordt binnen dit volume
afgelost: \cref{ch:b1:continuity} bewijst de \omterm{def:b1:logic:statement}{uitspraak} over de bijectie
(de stelling over de monotone inverse, steunend op de
tussenwaardestelling), en \cref{ch:b1:derivative} bewijst de regel
$(f^{-1})' = 1/(f' \circ f^{-1})$ die stilzwijgend elke formule van
\cref{prop:b1:functions:arcderiv} en \cref{prop:b1:functions:invhyp}
produceerde. Die hoofdstukken lezen met $\arcsin$ en
$\operatorname{arsinh}$ in gedachten --- als de uitgewerkte voorbeelden
waarvoor de theorie gebouwd is --- is de bedoelde weg rond de schijnbare
cirkelredenering.
\end{remark}

\section{Oefeningen}

\begin{exercise}[$\star$]\label{exo:b1:functions:1}
Bereken zonder rekenmachine: $\arcsin\frac{\sqrt 3}{2}$;
$\;\arccos\bigl(-\frac 12\bigr)$; $\;\arctan(-1)$;
$\;\arcsin\bigl(\sin\frac{5\pi}{6}\bigr)$;
$\;\arccos\bigl(\cos\frac{7\pi}{4}\bigr)$.
\end{exercise}

\begin{exercise}[$\star$]\label{exo:b1:functions:2}
Geef het definitiegebied van $f(x) = \arcsin(2x - 1)$ en bereken $f'$
waar die bestaat. Dezelfde vragen voor $g(x) = \arctan\sqrt{x}$.
\end{exercise}

\begin{exercise}[$\star$]\label{exo:b1:functions:3}
Bewijs dat voor alle $x \in \R$ geldt $\cosh x + \sinh x = \eu^x$, en dat
$(\cosh x + \sinh x)^n = \cosh nx + \sinh nx$ voor elke $n \in \Z$ (een
hyperbolische formule van de Moivre).
\end{exercise}

\begin{exercise}[$\star$]\label{exo:b1:functions:4}
Vereenvoudig $\cos(2\arcsin x)$ en $\sin(2\arctan x)$ tot algebraïsche
uitdrukkingen in $x$.
\end{exercise}

\begin{exercise}[$\star$]\label{exo:b1:functions:5}
Rangschik de functies $x^{100}$, $\;\eu^{\sqrt x}$, $\;(\ln x)^{1000}$,
$\;\eu^{x/100}$, $\;x^{\ln x}$ voor grote $x$ van traagst naar snelst
groeiend, met verantwoordingen gebaseerd op
\cref{prop:b1:functions:powerrules}.
\end{exercise}

\begin{exercise}[$\star\star$]\label{exo:b1:functions:6}
Bestudeer de functie $f(x) = \arctan\dfrac{2x}{1 - x^2}$ op haar
domein: bereken $f'$, vergelijk met $(2\arctan x)'$, en druk $f(x)$ op
elk van de drie intervallen van het domein uit in $\arctan x$.
\end{exercise}

\begin{exercise}[$\star\star$]\label{exo:b1:functions:7}
Los op in $\R$: $\;\cosh x = 2$; en vervolgens $5\cosh x - 4 \sinh x =
3$ (druk de oplossingen uit met logaritmen). \emph{Aanwijzing voor de
tweede: schrijf alles met $u = \eu^x$.}
\end{exercise}

\begin{exercise}[$\star\star$]\label{exo:b1:functions:8}
Bewijs de identiteit $\arctan\frac{1}{2} + \arctan\frac{1}{3} =
\frac{\pi}{4}$, en vervolgens de formule van
Machin\index{formule van Machin}
\[
4\arctan\frac 15 - \arctan\frac{1}{239} = \frac{\pi}{4}.
\]
\emph{Aanwijzing: bereken de tangens van beide leden met de somformule
en beheers in welk interval elk lid ligt.}
\end{exercise}

\begin{exercise}[$\star\star$]\label{exo:b1:functions:9}
Bewijs dat voor alle $x \geq 0$ geldt $\;x - \dfrac{x^3}{6} \leq \sin x
\leq x$ \emph{(bestudeer de opeenvolgende afgeleiden van de
verschillen)}, en leid af dat $\lim_{x \to 0^+} \frac{x - \sin x}{x^3} =
\frac 16$ aannemelijk is --- de limiet zelf wordt in
\cref{ch:b1:taylor} vastgelegd.
\end{exercise}

\begin{exercise}[$\star\star\star$]\label{exo:b1:functions:10}
Zet voor $x \in \intcc{-1}{1}$
$h(x) = \arcsin\bigl(2x\sqrt{1 - x^2}\,\bigr) - 2\arcsin x$. Je zou uit
de identiteit $\sin 2y = 2\sin y\cos y$ verwachten dat $h = 0$ --- maar
$h$ is \emph{niet} identiek nul. Bereken $h'$ op de open intervallen
waar die bestaat, evalueer $h$ in goedgekozen punten, en geef de
volledige stuksgewijs constante beschrijving van $h$ op $\intcc{-1}{1}$.
\end{exercise}

\begin{exercise}[$\star\star\star$]\label{exo:b1:functions:11}
(Gudermann-functie) Zij $g(x) = \arctan(\sinh x)$ voor $x \in \R$.
Bewijs dat $g$ een oneven, strikt stijgende bijectie van $\R$ op
$\intoo{-\frac\pi2}{\frac\pi2}$ is, dat $g'(x) = \frac{1}{\cosh x}$, en
dat $\tan g(x) = \sinh x$, $\;\sin g(x) = \tanh x$, $\;\cos g(x) =
\frac{1}{\cosh x}$: de functie $g$ verbindt de goniometrie met de
hyperbolische goniometrie zonder complexe getallen.
\end{exercise}

\begin{exercise}[$\star\star$]\label{exo:b1:functions:12}
Bewijs dat
\[
\arctan 1 + \arctan 2 + \arctan 3 = \pi .
\]
\emph{Aanwijzing: bereken eerst $\arctan 2 + \arctan 3$ en lokaliseer de
som als in \cref{ex:b1:functions:arctansum}; let erop dat de noemer van
de somformule hier negatief is.}
\end{exercise}

\section{Opgave: de vergelijking $x^y = y^x$}

\begin{problem}\label{pb:b1:functions:1}
Welke paren positieve getallen voldoen aan $x^y = y^x$? Iedereen kent
één toevallig ogend voorbeeld, $2^4 = 4^2 = 16$; deze opgave laat zien
dat daar niets toevalligs aan is. De hele vergelijking wordt geregeerd
door het verloop van de ene functie
\[
f(t) = \frac{\ln t}{t} \qquad (t > 0),
\]
waarvan de studie het volgende oplevert: de volledige
oplossingsverzameling (een diagonaal plus één gekromde tak door $(\eu,
\eu)$), een rationale parametrisering van die tak, het feit dat $(2, 4)$
haar \emph{enige} geheeltallige punt is, en de classificatie van al haar
rationale punten --- plus, als dividend, de vergelijking van $\eu^\pi$
met $\pi^\eu$ en de monotone convergentie van $\bigl(1 +
\frac1n\bigr)^n$ naar $\eu$.
\index{machtsfunctie}

\medskip
\noindent\textbf{Deel I --- De functie $f(t) = \ln t / t$.}
\begin{enumerate}
  \item Verantwoord dat $f$ afleidbaar is op $\intoo0{+\infty}$,
        bereken $f'$ en stel de verlooptabel op: $f$ stijgt strikt op
        $\intoc0\eu$, daalt strikt op $\intco\eu{+\infty}$, met maximum
        $f(\eu) = \frac1\eu$.
  \item Bepaal de limieten van $f$ in $0^+$ en $+\infty$
        (\cref{prop:b1:functions:powerrules}), het teken van $f$
        (negatief op $\intoo01$, nul in $1$, positief daarna), en
        schets de grafiek.
  \item Ga door rechtstreeks rekenen na dat $f(2) = f(4)$. (Houd die
        gelijkheid in gedachten: de hele opgave komt eruit voort.)
  \item Zij $c \in \R$. Bespreek, naar de waarde van $c$, het aantal
        oplossingen van $f(t) = c$: precies één voor $c \leq 0$;
        precies twee (één in $\intoo1\eu$, één in
        $\intoo\eu{+\infty}$) voor $0 < c < \frac1\eu$; precies één
        voor $c = \frac1\eu$; geen voor $c > \frac1\eu$.
  \item Leid af dat $\eu^x > x^\eu$ voor elke $x > 0$ met $x \neq \eu$,
        en beslis in het bijzonder welke van $\eu^\pi$ en $\pi^\eu$ de
        grootste is.
\end{enumerate}

\medskip
\noindent\textbf{Deel II --- De vergelijking en haar kromme.} In dit
deel is $x, y > 0$.
\begin{enumerate}[resume]
  \item Toon aan dat $x^y = y^x \iff f(x) = f(y)$.
  \item Leid de structuur van de oplossingsverzameling af: alle
        diagonale paren $(x, x)$; en de \emph{niet-triviale} paren ($x
        \neq y$), waarvoor geldt dat beide coördinaten $> 1$ zijn en
        dat de ene in $\intoo1\eu$ ligt terwijl de andere in
        $\intoo\eu{+\infty}$ ligt.
  \item Parametriseer de niet-triviale paren: toon aan dat, met $y =
        tx$ voor $t > 0$ en $t \neq 1$, de vergelijking $x^y = y^x$
        afdwingt dat
        \[
        x(t) = t^{\frac1{t-1}},
        \qquad
        y(t) = t\,x(t) = t^{\frac{t}{t-1}},
        \]
        en dat omgekeerd elk zo'n paar een oplossing is.
  \item Ga na dat $t = 2$ het paar $(2, 4)$ geeft, en bewijs de
        symmetrie $x(1/t) = y(t)$, $y(1/t) = x(t)$: de parameter
        omkeren verwisselt de twee coördinaten.
  \item Bepaal de limieten van $x(t)$ en $y(t)$ voor $t \to 1$ (beide
        gaan naar $\eu$), voor $t \to +\infty$ ($x \to 1$, $y \to
        +\infty$) en voor $t \to 0^+$ ($x \to +\infty$, $y \to 1$).
        Beschrijf de resulterende tak: een kromme met de rechten $x =
        1$ en $y = 1$ als asymptoten, die de diagonaal snijdt in $(\eu,
        \eu)$.
  \item Toon aan dat $t \mapsto x(t)$ strikt dalend is op
        $\intoo1{+\infty}$. \emph{(Bestudeer $h(t) = 1 - \frac1t - \ln
        t$, waarvan het teken de afgeleide van $\ln x(t) = \frac{\ln
        t}{t - 1}$ regelt.)}
  \item Besluit dat de niet-triviale oplossingen een strikt dalende
        bijectie $\varphi \colon \intoo1\eu \to \intoo\eu{+\infty}$
        definiëren met $x^{\varphi(x)} = \varphi(x)^x$, en dat $\varphi$
        een involutie van de tak is: $\varphi(\varphi(x)) = x$ overal
        waar beide leden gedefinieerd zijn.
\end{enumerate}

\medskip
\noindent\textbf{Deel III --- Geheeltallige en rationale punten.}
\begin{enumerate}[resume]
  \item Bewijs dat $(2, 4)$ en $(4, 2)$ de enige niet-triviale
        \emph{geheeltallige} oplossingen van $x^y = y^x$ zijn.
  \item Pas voor $n \in \N^*$ de parametrisering toe met $t = 1 +
        \frac1n$ en toon aan dat
        \[
        x_n = \Bigl(1 + \frac1n\Bigr)^{n},
        \qquad
        y_n = \Bigl(1 + \frac1n\Bigr)^{n+1}
        \]
        voor elke $n$ een niet-triviale \emph{rationale} oplossing is,
        met $(x_1, y_1) = (2, 4)$.
  \item Zij omgekeerd $(x, y)$ een niet-triviale rationale oplossing
        met $y > x$, en schrijf $t = y/x = 1 + \frac rs$
        onvereenvoudigbaar ($r, s \in \N^*$). Toon met $x = t^{s/r}$ en
        met de aangenomen uniciteit van de priemontbinding (vertrouwd
        van school; bewezen in \cref{ch:b1:arith}) aan dat $x$
        rationaal afdwingt dat zowel $s$ als $s + r$ een $r$-de macht
        van een geheel getal is.
  \item Toon met het binomium aan dat $a^r - b^r = r$ geen geheeltallige
        oplossingen $a > b \geq 1$ heeft wanneer $r \geq 2$, en
        besluit: de rationale oplossingen van $x^y = y^x$ zijn precies
        de paren $(x_n, y_n)$ van vraag 14 (en hun verwisselingen).
  \item Ga het geval $n = 2$ numeriek na: bereken $f(9/4)$ en $f(27/8)$
        tot op vier decimalen en controleer dat ze overeenstemmen.
\end{enumerate}

\medskip
\noindent\textbf{Deel IV --- Dividenden.}
\begin{enumerate}[resume]
  \item Bewijs met de vragen 7 en 11, zonder verder rekenwerk, dat de
        rij $x_n = \bigl(1 + \frac1n\bigr)^n$ strikt stijgend is met
        $x_n < \eu$, dat $y_n = \bigl(1 + \frac1n\bigr)^{n+1}$ strikt
        dalend is met $y_n > \eu$, en dat beide naar $\eu$ convergeren.
        (De parameter $t_n = 1 + \frac1n$ daalt naar $1$.)
  \item Stel de algemene vergelijkingsregel op voor $1 < a < b$: is
        $\eu \leq a$, dan $a^b > b^a$; is $b \leq \eu$, dan $a^b <
        b^a$; en toon met de twee voorbeelden $(2, 3)$ en $(2, 5)$ aan
        dat in het gemengde geval $a < \eu < b$ beide uitkomsten
        werkelijk voorkomen.
  \item Neem aan dat de involutie $\varphi$ van vraag 12 afleidbaar is
        in $\eu$ (dat is ze). Differentieer de identiteit
        $\varphi(\varphi(x)) = x$ in $x = \eu$ en leid af dat
        $\varphi'(\eu) = -1$: de tak snijdt de diagonaal loodrecht.
  \item Vind het unieke oplossingspaar met $y = 3x$ in gesloten vorm,
        en ga het numeriek na tot op vier decimalen met behulp van $f$.
  \item Bepaal onder de getallen $\sqrt2$, $\sqrt[3]3$, $\sqrt[4]4$,
        $\sqrt[5]5$ de grootste en wijs de twee aan die gelijk zijn.
        \emph{(Vergelijk $n^{1/n} = \eu^{f(n)}$.)}
\end{enumerate}

\medskip
\noindent\textbf{Deel V --- Synthese.}
\begin{enumerate}[resume]
  \item Beschrijf de volledige oplossingsverzameling van $x^y = y^x$ in
        het kwadrant $x, y > 0$ --- diagonaal plus tak, hun snijpunt
        $(\eu, \eu)$, de asymptoten, het geheeltallige punt $(2, 4)$,
        de rationale punten die zich ophopen bij $(\eu, \eu)$ --- in
        een vorm die je uit het hoofd zou kunnen schetsen.
  \item Waar precies gebruikte de opgave: (i) de \omterm{prop:b1:functions:powerrules}{groeivergelijkingen}
        van \cref{prop:b1:functions:powerrules}; (ii) de
        tussenwaardestelling (via de \omterm{def:b1:logic:statement}{uitspraken} over bijecties); (iii)
        de aangenomen uniciteit van de priemontbinding? Eén zin per
        onderdeel.
  \item Moraal, in een korte alinea: één verlooptabel loste een
        vergelijking in twee onbekenden op, classificeerde haar
        rationale punten en bewees de monotone convergentie van
        $\bigl(1 + \frac1n\bigr)^n$ --- becommentarieer die zuinigheid
        en noem waar elke draad later in het volume geïndustrialiseerd
        wordt (\cref{ch:b1:seq} voor de rij, \cref{ch:b1:derivative}
        voor verlooptabellen, \cref{ch:b1:taylor} voor de precisie die
        de tabel mist).
\end{enumerate}
\end{problem}
